% from ega0-20-fr.tex
\oldpage[0\textsubscript{IV-1}]{134}
(i) Comme $u$ est surjectif, on a $\operatorname{D\acute{e}r}_B(C,L)=0$ pour tout $C$-module
$L$ en vertu de (20.1.1). La suite exacte (20.2.3.1) devient donc
\begin{equation}
\tag{20.5.12.2}
\label{0.20.5.12.2-fr}
  0\longrightarrow\operatorname{D\acute{e}r}_A(C,L)\xrightarrow{u^0}\operatorname{D\acute{e}r}_A(B,L)
  \xrightarrow{\partial}\operatorname{Exalcom}_B(C,L)
  \xrightarrow{v^1}\operatorname{Exalcom}_A(C,L)
  \xrightarrow{u^1}\operatorname{Exalcom}_A(B,L),
\end{equation}
où $v$ est l'homomorphisme $A\to B$. Rappelons d'autre part (18.3.8) que
$\operatorname{Exalcom}_B(C,L)$ s'identifie canoniquement à
$\operatorname{Hom}_C(\mathfrak{R}/\mathfrak{R}^2,L)$; on déduit donc de (20.5.12.2)
et (20.4.8) la suite exacte
\begin{equation}
\tag{20.5.12.3}
\label{0.20.5.12.3-fr}
\begin{split}
  0\longrightarrow\operatorname{Hom}_C(\Omega^1_{C/A},L)
  \longrightarrow\operatorname{Hom}_C(\Omega^1_{B/A}\otimes_B C,L)
  \xrightarrow{\varphi}\operatorname{Hom}_C(\mathfrak{R}/\mathfrak{R}^2,L)\\
  \xrightarrow{\psi}\operatorname{Ker}(u^1)\longrightarrow 0
\end{split}
\end{equation}
avec $\varphi=\partial\circ\eta^{-1}$ et $\psi=v^1\circ\eta$. Revenant
aux définitions de $\partial$ (20.2.2) et de $\eta$ (18.3.8), on constate
aussitôt que $\varphi$ est précisément l'homomorphisme
$\operatorname{Hom}(\delta_{C/B/A},1_L)$. L'existence de la suite exacte formée des
4 premiers termes de (20.5.12.3) montre donc que la suite (20.5.12.1)
est exacte (Bourbaki, \emph{Alg.}, chap. II, 3\textsuperscript{e} éd.,
\S~2, n\textsuperscript{o}~1, th.~1).

\medskip
(ii) Appliquons à $B$ et à l'idéal $\mathfrak{R}^2$ la suite exacte
(20.5.12.1), ce qui donne
\begin{equation}
\tag{20.5.12.4}
\label{0.20.5.12.4-fr}
  \mathfrak{R}^2/\mathfrak{R}^4\longrightarrow
  \Omega^1_{B/A}\otimes_B E\longrightarrow\Omega^1_{E/A}\longrightarrow 0,
\end{equation}
d'où, en tensorisant par $C$ (considéré comme $E$-algèbre), la suite
exacte
\[
  \mathfrak{R}^2/\mathfrak{R}^4\otimes_E C\longrightarrow
  \Omega^1_{B/A}\otimes_B C\longrightarrow
  \Omega^1_{E/A}\otimes_E C\longrightarrow 0.
\]
Or, si $x,y$ sont deux éléments de $\mathfrak{R}$, et $\zeta$ la classe
de $xy$ mod. $\mathfrak{R}^4$, l'image
$\delta'(\zeta\otimes1)$ est par définition
$d_{B/A}(xy)\otimes1=(xd_{B/A}(y)+yd_{B/A}(x))\otimes1$; mais comme les
images de $x$ et de $y$ dans $C$ sont nulles, on a aussi
$d_{B/A}(xy)\otimes1=0$ dans $\Omega^1_{B/A}\otimes_B C$, ce qui prouve
notre assertion.

\medskip
(iii) Dire que $\delta_{C/B/A}$ est inversible à gauche revient, compte
tenu de l'exactitude de (20.5.12.1), à dire que la suite
\begin{equation}
\tag{20.5.12.5}
\label{0.20.5.12.5-fr}
  0\longrightarrow\mathfrak{R}/\mathfrak{R}^2
  \xrightarrow{\delta_{C/B/A}}\Omega^1_{B/A}\otimes_B C
  \xrightarrow{u_{C/B/A}}\Omega^1_{C/A}\longrightarrow0
\end{equation}
est \emph{exacte et scindée}, et il revient au même (Bourbaki,
\emph{loc. cit.}) de dire que $\operatorname{Ker}(u^1)=0$ dans la suite exacte
(20.5.12.3) pour tout $L$, ce qui montre l'équivalence des conditions
a) et b) (cf. (18.3.6.2)).

Le fait que b) entraîne c) provient de ce que l'image réciproque par
$u:B\to C$ de la $A$-extension $E$ de $C$ par
$\mathfrak{R}/\mathfrak{R}^2$ est $A$-triviale, $u$ étant composée des
homomorphismes canoniques $B\to E\to C$ (18.1.6). Inversement, c)
entraîne b), car toute $B$-extension de $C$ par $L$ est $B$-équivalente
à $E\oplus_{\mathfrak{R}/\mathfrak{R}^2}L$ (18.3.8), autrement dit sa
classe est l'image de la classe de $E$ (considéré comme $B$-extension)
par l'homomorphisme

